\section{Практическая часть}

\subsection{Описание лабораторной установки}

Лабораторная установка реализована в виде приложения \texttt{Aver\_estimation.exe}, созданного средствами MATLAB Compiler. По методическому пособию запуск приложения требует установленной среды MATLAB Compiler Runtime соответствующей версии.

Установка состоит из четырёх функциональных блоков:

\begin{itemize}
	\item генератор сигналов --- формирует дискретный гауссовский шум и позволяет выбирать $\tau_\text{корр} = 10$, $30$ или $100$;
	\item усреднитель --- задаёт число отсчётов $N$, число реализаций $M$, временной интервал $\Delta t$ и доверительную вероятность $\beta$;
	\item анализатор --- отображает реализацию, усреднённую реализацию, СПМ исходного и усреднённого процесса, а также гистограмму оценок;
	\item блок оценок --- показывает среднее, СКО и доверительный интервал.
\end{itemize}

В эксперименте использовались кнопки «Run», «Усреднение», «СКО(N)» и «СКО(delta\_t)». Кроме того, применялась кнопка «PDF». Кнопка «Reset» использовалась для очистки области графиков перед началом новой серии, когда это требовалось методическим пособием.

Ниже приведены скриншоты из пользовательского архива «Для отчёта.rar» и три повторно проверенных снимка задания~5, переданных отдельно. Для серийных экспериментов используются итоговые изображения, на которых одновременно видны все кривые серии. Промежуточные снимки, содержащие только одну или две кривые, не дублируются в основном тексте.

\subsection{Задание 1. Реализации случайного процесса и его спектры}

Исследованы реализации гауссовского шума и их спектральные плотности мощности при $\tau_\text{корр} = 10$, $30$ и $100$. Для каждого значения времени корреляции получена реализация сигнала и график СПМ.

\begin{figure}[H]
	\centering
	\begin{minipage}{0.40\textwidth}
		\centering
		\includegraphics[width=\linewidth]{fig01}
		\subcaption{Реализация сигнала, $\tau_\text{корр} = 10$}
		\label{fig:pr:1a}
	\end{minipage}\hfill
	\begin{minipage}{0.40\textwidth}
		\centering
		\includegraphics[width=\linewidth]{fig02}
		\subcaption{Реализация сигнала, $\tau_\text{корр} = 30$}
		\label{fig:pr:1b}
	\end{minipage}

	\vspace{0.5em}
	\begin{minipage}{0.40\textwidth}
		\centering
		\includegraphics[width=\linewidth]{fig03}
		\subcaption{Реализация сигнала, $\tau_\text{корр} = 100$}
		\label{fig:pr:1c}
	\end{minipage}\hfill
	\begin{minipage}{0.40\textwidth}
		\centering
		\includegraphics[width=\linewidth]{fig04}
		\subcaption{СПМ, $\tau_\text{корр} = 10$}
		\label{fig:pr:1d}
	\end{minipage}

	\vspace{0.5em}
	\begin{minipage}{0.40\textwidth}
		\centering
		\includegraphics[width=\linewidth]{fig05}
		\subcaption{СПМ, $\tau_\text{корр} = 30$}
		\label{fig:pr:1e}
	\end{minipage}\hfill
	\begin{minipage}{0.40\textwidth}
		\centering
		\includegraphics[width=\linewidth]{fig06}
		\subcaption{СПМ, $\tau_\text{корр} = 100$}
		\label{fig:pr:1f}
	\end{minipage}

	\caption{Реализации сигнала и спектральные плотности мощности при различных временах корреляции}
	\label{fig:pr:1}
\end{figure}

При $\tau_\text{корр} = 10$ сигнал меняется наиболее быстро; при $\tau_\text{корр} = 30$ колебания становятся более плавными; при $\tau_\text{корр} = 100$ наблюдаются наиболее протяжённые во времени участки. По спектрам видно, что увеличение времени корреляции сопровождается сужением характерной полосы спектральных составляющих около нулевой частоты. Качественно это соответствует обратной связи между характерной шириной спектра и временем корреляции: $\Delta f \sim 1/\tau_\text{корр}$.

\subsection{Задание 2. Поведение оценки среднего при изменении числа отсчётов}

На первом этапе при $\tau_\text{корр} = 10$, $M = 256$, $\Delta t = 1$ и $N = 1$ выполнено усреднение ансамбля. В блоке оценок получены среднее значение $5{,}79564$ и СКО $14{,}525$. Соответствующая оценка дисперсии равна:
\[
	D[x] \approx \sigma^2 = 14{,}525^2 = 210{,}98 .
\]

\begin{figure}[H]
	\centering
	\begin{minipage}{0.47\textwidth}
		\centering
		\includegraphics[width=\linewidth]{fig07}
		\subcaption{Результат усреднения: $N = 1$, $M = 256$}
		\label{fig:pr:2a}
	\end{minipage}\hfill
	\begin{minipage}{0.47\textwidth}
		\centering
		\includegraphics[width=\linewidth]{fig08}
		\subcaption{Изменение оценок среднего при $N$ от $1$ до $128$ ($M = 8$)}
		\label{fig:pr:2b}
	\end{minipage}
	\caption{Результаты усреднения ансамбля и изменение оценок среднего}
	\label{fig:pr:2}
\end{figure}

На втором этапе использованы $M = 8$, $\Delta t = 1$ и максимальное $N = 128$. На графике показаны восемь траекторий оценки среднего при последовательном увеличении числа отсчётов. В начале разброс между реализациями велик, затем оценки группируются около общего значения.

Вертикальный размах по сохранённому графику можно оценить только приблизительно, поскольку линии частично перекрываются: при $n = 1$ он составляет около $40$--$45$ условных единиц; при $n = 40$ --- около $6$; при $n = 128$ --- около $3$--$4$. Ориентировочное отношение размахов составляет примерно $43 : 6 : 3{,}5$, или $1 : 0{,}14 : 0{,}08$ относительно $n = 1$. Это визуальная оценка по графику, а не точное вычисление выборочного СКО по исходным числовым данным.

Вертикальный разброс характеризует случайную изменчивость оценки среднего между независимыми реализациями при фиксированном числе усредняемых отсчётов. Его уменьшение при росте $n$ согласуется с законом больших чисел и с уменьшением дисперсии среднего.

\subsection{Задание 3. Зависимость СКО оценки от числа усредняемых отсчётов}

Получены две серии зависимостей $\sigma(N)$ при $M = 256$, $N = 64$ и трёх значениях $\Delta t = 1$, $4$, $12$. Первая серия соответствует $\tau_\text{корр} = 10$, вторая --- $\tau_\text{корр} = 100$. На итоговых графиках каждая серия содержит все три кривые (\cref{fig:pr:3}).

\begin{figure}[H]
	\centering
	\begin{minipage}{0.47\textwidth}
		\centering
		\includegraphics[width=\linewidth]{fig09}
		\subcaption{Итоговая серия при $\tau_\text{корр} = 10$: $\Delta t = 1, 4, 12$}
		\label{fig:pr:3a}
	\end{minipage}\hfill
	\begin{minipage}{0.47\textwidth}
		\centering
		\includegraphics[width=\linewidth]{fig10}
		\subcaption{Итоговая серия при $\tau_\text{корр} = 100$: $\Delta t = 1, 4, 12$}
		\label{fig:pr:3b}
	\end{minipage}
	\caption{Зависимости СКО оценки от числа усредняемых отсчётов}
	\label{fig:pr:3}
\end{figure}

С увеличением $N$ СКО уменьшается. При фиксированном $\tau_\text{корр}$ увеличение $\Delta t$ ослабляет корреляцию соседних измерений и снижает погрешность. При $N = 64$ для $\tau_\text{корр} = 10$ по графику получаются ориентировочные СКО около $4{,}1$; $2{,}0$ и $1{,}7$ при $\Delta t = 1$, $4$ и $12$; для $\tau_\text{корр} = 100$ --- около $12{,}2$; $6{,}0$ и $1{,}9$.

В перечне вопросов к заданию~3 методичка просит объяснить уменьшение СКО при фиксированном $N$ с ростом времени корреляции. Однако по сохранённым графикам зависимости $\sigma(N)$ при $\tau_\text{корр} = 100$ кривые расположены выше, чем при $\tau_\text{корр} = 10$ (особенно для $\Delta t = 1$ и $4$). Следовательно, наблюдаемое направление не подтверждает формулировку вопроса методички. По теоретической формуле дисперсии оценки корреляция отсчётов увеличивает погрешность; поэтому увеличение времени корреляции при фиксированном $\Delta t$ обычно означает более сильную зависимость соседних отсчётов. В отчёте сохранены фактически наблюдённые графики, а несоответствие с формулировкой методички обозначено явно. При $N = 1$ усреднения нет, поэтому СКО оценки определяется СКО исходного процесса. Нулевой начальный наклон при $N = 1$ связывается с тем, что ковариационная функция стационарного процесса имеет экстремум при нулевом временном сдвиге.

\subsection{Задание 4. Зависимость СКО оценки от временного интервала}

Исследована зависимость СКО от $\Delta t$ в диапазоне $1$--$32$ для $\tau_\text{корр} = 10$, $30$, $100$. Получены две серии: при $N = 4$ и при $N = 32$; в обеих сериях $M = 256$ (\cref{fig:pr:4}).

\begin{figure}[H]
	\centering
	\begin{minipage}{0.47\textwidth}
		\centering
		\includegraphics[width=\linewidth]{fig11}
		\subcaption{Итоговая серия при $N = 4$: $\tau_\text{корр} = 10, 30, 100$}
		\label{fig:pr:4a}
	\end{minipage}\hfill
	\begin{minipage}{0.47\textwidth}
		\centering
		\includegraphics[width=\linewidth]{fig12}
		\subcaption{Итоговая серия при $N = 32$: $\tau_\text{корр} = 10, 30, 100$}
		\label{fig:pr:4b}
	\end{minipage}
	\caption{Зависимости СКО оценки от интервала между отсчётами}
	\label{fig:pr:4}
\end{figure}

При $\Delta t$, близком к нулю, соседние отсчёты почти совпадают по статистическим свойствам, поэтому небольшое изменение интервала слабо меняет их корреляцию. Поскольку ковариационная функция стационарного процесса чётная, её производная в нуле равна нулю; этим объясняется нулевой начальный наклон зависимости $\sigma(\Delta t)$.

С ростом $\Delta t$ корреляция между отсчётами уменьшается, и СКО оценки падает. При достаточно больших интервалах отсчёты становятся практически независимыми, и кривая приближается к пределу $\sigma_x/\sqrt{N}$. Поэтому, чтобы уменьшить этот предел в $2$ раза, необходимо увеличить $N$ в $4$ раза (например, с $4$ до $16$ или с $32$ до $128$). При фактически исследованном переходе $N = 4 \to 32$ теоретический коэффициент уменьшения равен $\sqrt{32/4} = \sqrt{8} \approx 2{,}83$.

Время корреляции оценивается по характерному «излому» кривой --- области, после которой дальнейшее увеличение $\Delta t$ даёт всё меньший выигрыш. Для $\tau_\text{корр} = 10$ этот переход заметен примерно при $\Delta t$ порядка $10$; для $\tau_\text{корр} = 30$ --- в области нескольких десятков; для $\tau_\text{корр} = 100$ диапазона $\Delta t \leqslant 32$ недостаточно, чтобы уверенно увидеть полный выход на плато.

\subsection{Задание 5. Определение характеристик по спектральной плотности мощности}

\subsubsection{Исходный процесс (задание 5.1)}

При $\tau_\text{корр} = 10$ получен спектр исходного процесса в режиме «СПМ» при $N = 1$ и $M = 1$. Согласно методическому пособию, мощность в нулевой частотной полосе пропорциональна квадрату среднего значения. Если $A$ --- значение СПМ в нулевой частотной полосе, а $\Delta \omega = 1/2048$, то:
\[
	\langle x \rangle^2 \approx A \cdot \frac{1}{2048}, \qquad |\langle x \rangle| \approx \sqrt{\frac{A}{2048}} .
\]

По высоте нулевого пика на повторно проверенном скриншоте можно снять только приближённое значение $A \approx 1{,}2 \cdot 10^4$. Тогда графическая оценка модуля среднего составляет $|\langle x \rangle| \approx \sqrt{1{,}2 \cdot 10^4 / 2048} \approx 2{,}4$. Это именно оценка по растровому изображению, а не точное значение из числового массива спектра; знак среднего по СПМ таким способом не определяется.

Дисперсия исходного процесса определяется площадью под непрерывной частью СПМ; в методическом пособии она также описывается как произведение эффективной ширины спектра на эффективное значение СПМ: $D[x] \approx \Delta \omega_\text{эфф} \cdot S_\text{эфф}$. В переданном файле имеется график, но нет массива спектральных отсчётов и курсорных значений границ эффективной полосы. Поэтому точное численное значение дисперсии по данному изображению восстановить надёжно нельзя; в отчёте не подменяем его полями блока оценок.

\begin{figure}[H]
	\centering
	\includegraphics[width=0.6\linewidth]{fig13}
	\caption{СПМ исходного процесса при $\tau_\text{корр} = 10$, $N = 1$, $M = 1$}
	\label{fig:pr:5}
\end{figure}

\textbf{Сопоставление с заданием~2.} В задании~2 по ансамблю получены среднее $5{,}79564$ и СКО $14{,}525$ (дисперсия около $210{,}98$). Графическое считывание пика СПМ даёт лишь ориентировочный модуль среднего около $2{,}4$. Эти оценки заметно различаются; ввиду ограниченной точности считывания пика по изображению и отсутствия исходного массива спектра данное расхождение фиксируется, а не устраняется подбором значения. Числа $5{,}79564$ и $14{,}525$ в правом блоке обновлённого скриншота~5.1 сохранены интерфейсом после предыдущего эксперимента и не используются как отдельный расчёт по СПМ.

\subsubsection{Усреднённый процесс (задание 5.2)}

Для усреднённого процесса при $\tau_\text{корр} = 10$ и $\Delta t = 1$ получены спектры при $M = 2$ для $N = 4$ и $N = 32$. В соответствии с методичкой временное усреднение можно рассматривать как линейный фильтр: постоянная составляющая сохраняется, а эффективная полоса пропускания около нулевой частоты с увеличением времени усреднения сужается. Поэтому при переходе от $N = 4$ к $N = 32$ непрерывная часть спектра становится более узкой, а дисперсия оценки уменьшается (\cref{fig:pr:6}).

\begin{figure}[H]
	\centering
	\begin{minipage}{0.47\textwidth}
		\centering
		\includegraphics[width=\linewidth]{fig14}
		\subcaption{Усреднённая СПМ при $\tau_\text{корр} = 10$, $N = 4$, $M = 2$}
		\label{fig:pr:6a}
	\end{minipage}\hfill
	\begin{minipage}{0.47\textwidth}
		\centering
		\includegraphics[width=\linewidth]{fig15}
		\subcaption{Усреднённая СПМ при $\tau_\text{корр} = 10$, $N = 32$, $M = 2$}
		\label{fig:pr:6b}
	\end{minipage}
	\caption{Спектральные плотности мощности усреднённого процесса}
	\label{fig:pr:6}
\end{figure}

На обоих графиках высота пика около нулевой частоты остаётся одного порядка, что согласуется с сохранением постоянной составляющей при усреднении. При $N = 32$ непрерывная часть заметно сильнее сосредоточена около нулевой частоты и занимает меньшую полосу, чем при $N = 4$; это соответствует уменьшению дисперсии оценки. Методичка предписывает определить среднее и дисперсию по этим спектрам и сравнить их с заданием~2, однако в доступных материалах представлены только растровые изображения без точных массивов спектральных значений. Поэтому точные интегральные значения дисперсии для $N = 4$ и $N = 32$ не заявляются. Числа в правом блоке «Оценки» на скриншотах не использованы как свежие спектральные оценки: интерфейс сохранял значения, полученные в предшествующих опытах.

\subsection{Задание 6.1. Гистограмма оценок и доверительный интервал}

Исследовано изменение гистограммы оценок среднего при $\tau_\text{корр} = 10$, $M = 256$, $\Delta t = 1$ и $\beta = 0{,}95$. Эксперименты выполнены для $N = 1$, $4$ и $32$ (\cref{fig:pr:7}).

\begin{figure}[H]
	\centering
	\begin{minipage}{0.47\textwidth}
		\centering
		\includegraphics[width=\linewidth]{fig16}
		\subcaption{Гистограмма и доверительный интервал, $N = 1$}
		\label{fig:pr:7a}
	\end{minipage}\hfill
	\begin{minipage}{0.47\textwidth}
		\centering
		\includegraphics[width=\linewidth]{fig17}
		\subcaption{Гистограмма и доверительный интервал, $N = 4$}
		\label{fig:pr:7b}
	\end{minipage}

	\vspace{0.5em}
	\begin{minipage}{0.47\textwidth}
		\centering
		\includegraphics[width=\linewidth]{fig18}
		\subcaption{Гистограмма и доверительный интервал, $N = 32$}
		\label{fig:pr:7c}
	\end{minipage}
	\caption{Гистограммы оценок и доверительные интервалы}
	\label{fig:pr:7}
\end{figure}

По мере увеличения $N$ гистограмма становится более узкой и концентрируется вблизи центрального значения. Одновременно уменьшается доверительный интервал, что свидетельствует о повышении точности оценки среднего.

\begin{table}[H]
	\centering
	\caption{Доверительный интервал, полученный в задании 6.1}
	\label{tab:pr:1}
	\begin{tabular}{cccccc}
		\toprule
		$N$ & $\tau_\text{корр}$ & $M$ & $\Delta t$ & $\beta$ & $I_\beta$ по программе \\
		\midrule
		1  & 10 & 256 & 1 & 0,95 & 1,51876 \\
		4  & 10 & 256 & 1 & 0,95 & 0,769429 \\
		32 & 10 & 256 & 1 & 0,95 & 0,228206 \\
		\bottomrule
	\end{tabular}
\end{table}

\subsection{Задание 6.2. Зависимость доверительного интервала от $\beta$ и $N$}

При $\tau_\text{корр} = 10$, $M = 256$ и $\Delta t = 1$ исследована зависимость доверительного интервала от доверительной вероятности $\beta = 0{,}80$; $0{,}95$; $0{,}98$ для $N = 1$, $4$ и $32$. Значения $I_\beta$ считывались из блока оценок программы после каждого нажатия «PDF».

\begin{table}[H]
	\centering
	\caption{Зависимость доверительного интервала от доверительной вероятности}
	\label{tab:pr:2}
	\begin{tabular}{cccc}
		\toprule
		$N$ & $I_\beta$ при $\beta = 0{,}80$ & $I_\beta$ при $\beta = 0{,}95$ & $I_\beta$ при $\beta = 0{,}98$ \\
		\midrule
		1  & 1,300170 & 1,618780 & 2,188120 \\
		4  & 0,607883 & 0,716833 & 1,074840 \\
		32 & 0,166798 & 0,203898 & 0,295513 \\
		\bottomrule
	\end{tabular}
\end{table}

Для всех трёх $N$ доверительный интервал возрастает при увеличении $\beta$. При фиксированной $\beta$ интервал уменьшается с увеличением $N$. На \cref{fig:pr:8} приведена зависимость, построенная по девяти значениям, снятым с пользовательских скриншотов.

\begin{figure}[H]
	\centering
	\includegraphics[width=0.42\linewidth]{fig19}
	\caption{Зависимость доверительного интервала $I_\beta$ от доверительной вероятности $\beta$ по результатам задания 6.2}
	\label{fig:pr:8}
\end{figure}

\textbf{Примечание.} Значения $I_\beta$ при $N = 1$ и $\beta = 0{,}95$ в заданиях~6.1 и~6.2 различаются ($1{,}51876$ и $1{,}61878$), поскольку получены в отдельных запусках программы. Методичка рекомендует при малых $N$ (примерно до $15$--$20$), если дисперсия оценивается по той же выборке, учитывать распределение Стьюдента; в задании~6.2 программа использует гауссово приближение.
